% Pursuing Stacks (cotes 134-2 to 134-8): archive cote 134-6, site page 50 (PDF page 51), typescript page 368.
% typescript. Draft compared: ps5.tex:3836-3916.
% AI-transcribed by Claude Opus 5.5 (claude-opus-5-5, adaptive thinking, effort high), Message Batches API, batch ps3 (msgbatch_01XdFmuVVBAcgc7EFxdtCqsj), 2026-11-12; not yet checked by eye.
\scanpage{134-6/50}
\mspage{368}
general object in the family (5) \ $b_i^n$ rather than $b_i$ ). \ Thus the n.th
component of the cochain complex of the second member of (2) can be
explicitly written as
\[ (6) \qquad \mathrm{Hom}^n(L_\cdot^B,F) = \mathrm{Hom}(L_n^B,F) \simeq \hand{\bigoplus}_{i\in I_n} F(b_i) \quad , \]
and the \inserted{co}boundary operators between these components can be made explicit
in a similar way, \ by means of (possibly infinite) matrices, \ whose entries
% NOTE: the right edge of the scan cuts off the end of the next line
are $\underline{Z}$-linear combinations of maps from some $b_i^n$ to some $b_j^{n-1}$ \ $(i \hand{\in} I_n, j\hand{\in} I_{n-1}$.
We feel a little happier still when the direct sums (4) yielding the compo-%
nents $L_n$ are finite, \ i.e. the sets $I_n$ are finite, \ which also means that
$L_\cdot^B$ \ can be interpreted as a chain complex in the additive envelope Add(B)
of B, \ as contemplated in the first place \ - \ in which case the integrator
will be called \uline{``special''}.

The formula (2) immediately generalizes when F is replaced by a \xout{compl}
complex of \inserted{pre}sheaves $F^\bullet$ \ , \ with degrees \ bounded from below (NB as the nota-%
tion indicates, \ the differential operator is of degree +1), \ to
\[ (\over{\unsure{7'}}{7}) \qquad R\Gamma_B(F^\bullet) \xrightarrow{\sim} \mathrm{Hom}^{\bullet\bullet}(L_\cdot^B,F^\bullet) \quad , \]
where now the \xout{first} left hand side designates hypercohomology of B (\xout{or} i.e.
of the corresponding topos), \ viewed as an object of the right derived
category \ $D^+$(Ab) \ of the category of abelian groups, \ and where $\mathrm{Hom}^{\bullet\bullet}$
designates the double complex obtained by taking Hom's componentwise,
or more accurately, \ the object in $D^+$(Ab) \ defined by the associated simple
complex.

An interesting special case of (2') is obtained when starting with
a complex of abelian groups $K^\bullet$ \ bounded from below i.e. defining an object
of the right derived category $D^+$(Ab), \ and taking
\[ (8) \qquad F^\bullet = K_B^\bullet \ = \over{\illegible{1}}{p}{}_B^*(K^\bullet) \quad , \]
the corresponding \uline{constant} \uline{complex of presheaves} on B, \ which may be vie-%
wed equally as the inverse image of $K^\bullet$ \ by the projection
\[ (9) \qquad p_B \ : \ B \hand{\longrightarrow} \hand{\Delta}_o \ \text{(the final category)} \quad , \]
which geometrically interpretes as the canonical morphism of the topos
associated to \xout{the} B \ to the final (or ``one point'' ) topos. \ \unsure{The} second
member of (7) \ can be rewritten componentwise, \ using the adjunction
formula for the pair ($p_!^{ab}$ \ , $p^*$ ) (where the qualifying B is omitted now in
the notation p) :
\[ \mathrm{Hom}(L_n,p^*(K^m)) \simeq \mathrm{Hom}(p_!^{ab}(L_n),K^m) \quad , \]
so that (7) can be rewritten \ as

\margin{50}
